\begin{corollary}[2.3.8]
\label{IV.2.3.8-fr}
\emph{(Zariski).} Sean $A$, $B$ dos anillos locales noetherianos,
$\mathfrak m$, $\mathfrak n$ sus ideales maximales respectivos,
$\varphi:A\to B$ un homomorfismo local. Se supone que se satisfacen las
hipótesis siguientes:
\begin{enumerate}
  \item[$1^\circ$] $B$ es una $A$-álgebra esencialmente de tipo finito
  \textup{(1.3.8)}.
  \item[$2^\circ$] El completado $\widehat A$ de $A$ para la topología
  $\mathfrak m$-ádica es íntegro.
  \item[$3^\circ$] $\varphi$ es inyectivo.
\end{enumerate}
Entonces la topología $\mathfrak m$-ádica de $A$ está inducida por la
topología $\mathfrak n$-ádica de $B$.
\end{corollary}

Pongamos $B'=B\otimes_A\widehat A$; en virtud de $1^\circ$, $B'$ es de la
forma $S^{-1}(C\otimes_A\widehat A)$, donde $C$ es una $A$-álgebra de tipo
finito y $S$ es una parte multiplicativa de $C$; por tanto, $B'$ es un anillo
\emph{noetheriano}. Como $A$ se identifica con un subanillo de $\widehat A$
$(\mathbf{0}_{\mathrm I},7.3.5)$, $A$ es íntegro por $2^\circ$. La hipótesis
$3^\circ$ implica entonces que existe un ideal primo $\mathfrak q$ de $B$
que induce el ideal 0 de $A$ $(\mathbf{0}_{\mathrm I},1.5.8)$ y, por consiguiente, el
homomorfismo local $A\to B/\mathfrak q$ es inyectivo. Es posible, pues,
limitarse a demostrar la conclusión de (2.3.8) añadiendo la hipótesis de que
$B$ es un anillo local \emph{íntegro}. Apliquemos (2.3.7, (ii)) a
$Y=\operatorname{Spec}(A)$, $X=\operatorname{Spec}(B)$,
$Y'=\operatorname{Spec}(\widehat A)$ y $X'=\operatorname{Spec}(B')$; como el
morfismo $Y'\to Y$ es plano y $X$
\oldpage[IV-2]{17}
(que es íntegro) domina el esquema íntegro
$Y$, toda componente irreducible de $X'$ domina el esquema (íntegro por
hipótesis) $Y'$. Si $y$, $x$, $y'$ son los puntos cerrados de $Y$, $X$ e
$Y'$, respectivamente, existe un punto $x'\in X'$ (por lo demás único) por
encima de $x$ e $y'$ (\textbf I, 3.4.9), y
$\operatorname{Spec}(\mathcal O_{x'})$ domina, por tanto,
$\operatorname{Spec}(\mathcal O_{y'})$; se tiene, en consecuencia, un
diagrama conmutativo de homomorfismos \emph{locales} de anillos locales noetherianos
\[
\xymatrix{
  B=\mathcal O_x \ar[r] & \mathcal O_{x'}\\
  A=\mathcal O_y \ar[u]^{\varphi} \ar[r]_{u} &
  \mathcal O_{y'}=\widehat A \ar[u]_{v}
}
\]
en el que $u$ y $v$ son \emph{inyectivos} (\textbf I, 1.2.7); identificando $A$ y
$\widehat A$ con subanillos de $\mathcal O_{x'}$, y designando por
$\mathfrak r$ el ideal maximal de $\mathcal O_{x'}$, la intersección de los
ideales $\mathfrak r^k\cap\widehat A$ es, por tanto, nula
$(\mathbf{0}_{\mathrm I},7.3.5)$; como $\widehat A$ es \emph{completo} y estos ideales son
abiertos en $\widehat A$, ello implica (Bourbaki, \emph{Alg. comm.},
cap.~III, \S~2, núm.~7, prop.~8) que la topología de $\widehat A$ está
\emph{inducida} por la topología $\mathfrak r$-preádica de $\mathcal O_{x'}$; \emph{a
fortiori}, lo mismo ocurre con la topología de $A$
$(\mathbf{0}_{\mathrm I},7.3.5)$. Por otra parte, se tiene
$\mathfrak n^k\cap A\subset\mathfrak r^k\cap A$, de modo que la topología
$\mathfrak n$-preádica de $B$ induce en $A$ una topología \emph{más fina} que la
topología $\mathfrak m$-preádica; pero, como
$\mathfrak m^k\subset\mathfrak n^k\cap A$, estas dos topologías son
\emph{idénticas}. Q.E.D.

\begin{remark}[2.3.9]
\label{IV.2.3.9-fr}
Veremos más adelante (7.8.3, (vii)) que, para los anillos locales noetherianos
$A$ más usuales en Geometría algebraica, la hipótesis de que $A$ es íntegro e
integralmente cerrado implica que lo mismo ocurre con $\widehat A$. Por ello,
en Geometría algebraica sobre un cuerpo base, se enuncia generalmente (2.3.8)
bajo la hipótesis de que $A$ es \emph{íntegro e integralmente cerrado}.
\end{remark}
